\section{First Principles}
\label{sec:first-principles}

\subsection{The Structure of the Argument}

The argument for the \dia{} rests on four premises and one logical step.
The premises are not contested.
Each has been endorsed, explicitly or implicitly, by both parties, by
Congress, and by the Supreme Court.
The logical step is elementary.
The conclusion is that \hh{} apportionment, applied within states via \ar{},
is the constitutionally correct method for congressional redistricting.

The argument may be stated formally as follows.

\begin{quotation}
\noindent
\textbf{Premise~1} (\wesberry{}, 1964): Congressional districts must have
equal population.\\[0.5em]
\textbf{Premise~2} (2~U.S.C.~\usca{}~2a(a), enacted 1941, upheld 1992):
The correct method for the interstate apportionment problem---allocating
congressional seats among states in proportion to population---is the
Huntington-Hill equal-proportions algorithm.\\[0.5em]
\textbf{Premise~3} (both parties, on record in courts across six decades):
The map should not be drawn by politicians with a personal stake in the
outcome.\\[0.5em]
\textbf{Premise~4} (logic): The interstate problem and the intrastate
problem are structurally isomorphic: both require apportioning $k$
indivisible units among sub-units in proportion to population,
minimizing percentage deviation from ideal allocation.\\[0.5em]
\textbf{Conclusion}: \hh{} should apply to the intrastate problem.
The algorithm that implements this is \ar{}.
The statute that mandates it is the \dia{}.
\end{quotation}

We develop each premise in turn.

\subsection{Premise 1: One Person, One Vote}

\wesberry{} v.\ \textit{Sanders}, 376 U.S. 1 (1964), is the foundational
principle of congressional redistricting.
Writing for the Court, Justice Black held that Article~I, \S2---``chosen by
the People of the several States''---means that ``as nearly as is practicable
one man's vote in a congressional election is to be worth as much as
another's.''\footnote{%
  \wesberry{}, 376 U.S. at 7--8, 18.
}
The Court declined to specify a precise tolerance, but the requirement was
clear: districts must be as equal in population as is practicable.

Nineteen years later, \karcher{} v.\ \textit{Daggett}, 462 U.S. 725 (1983),
converted this principle into a near-zero tolerance standard.
Any population deviation in a congressional district plan must be justified
by a legitimate state objective, and the state bears the burden of proving
the deviation was necessary to achieve that objective.\footnote{%
  \karcher{}, 462 U.S. at 730--731.
  The Court struck down a New Jersey plan with a maximum deviation of 0.6984\%
  because the state had not demonstrated that the deviation was required to
  achieve any legitimate objective.
}

The practical effect of \wesberry{} and \karcher{} together is that population
equality is, effectively, a hard constraint.
The target population for each district is $P/k$, where $P$ is the state's
total population and $k$ is its congressional seat count.
Each district must contain approximately $P/k$ persons.
This is not negotiable; it is the constitutional floor.

The \dia{} treats this as a hard constraint.
The parameter \texttt{ubvec[0]=1.001} in the \ar{} implementation
(the METIS population balance tolerance) imposes a $\pm 0.05\%$ deviation
ceiling---an order of magnitude tighter than the tightest maps that have
survived judicial review.
Premise~1 is satisfied by design.

\subsection{Premise 2: Huntington-Hill Is the Agreed Interstate Standard}

The Huntington-Hill method is codified at 2~U.S.C.~\usca{}~2a(a), which
provides that following each decennial census, the Clerk of the House of
Representatives shall ``determine the number of Representatives to which each
State is entitled'' using the ``method of equal proportions.''\footnote{%
  The statute does not name the method; the phrase ``method of equal
  proportions'' is the standard mathematical name for the algorithm
  devised by Joseph A.\ Hill and developed by Edward V.\ Huntington.
  See \textit{Department of Commerce v.\ Montana}, 503 U.S. 442, 448 (1992).
}

The method works as follows.
Assign one seat to each state (the constitutional minimum).
Then, iteratively, assign the next seat to the state with the highest
\emph{priority value}, defined as:
\[
  P_s(n) = \frac{\text{population}_s}{\sqrt{n(n+1)}}
\]
where $n$ is the state's current seat count and the divisor
$\sqrt{n(n+1)}$ is the geometric mean of $n$ and $n+1$.
Repeat until all 435 seats are assigned.

The geometric mean divisor minimizes the percentage difference between
any two states' average district populations.
This is the ``equal proportions'' property: Huntington-Hill is the method
that makes the seat allocation as proportional to population as possible,
in the sense of minimizing maximum percentage deviation.

The Supreme Court addressed Huntington-Hill directly in
\textit{United States Department of Commerce v.\ Montana},
503 U.S. 442 (1992), where Montana challenged an apportionment that
left it with one seat despite a population that would have entitled it to
nearly two seats under alternative methods.
The Court upheld the method, noting that the Constitution ``vests Congress
with the responsibility to configure and design the apportionment process,''
and that the equal-proportions method represents a ``rational decision'' within
Congress's discretion.\footnote{%
  \textit{Department of Commerce v.\ Montana}, 503 U.S. at 462--463.
}

No party has challenged the use of Huntington-Hill for the interstate
apportionment since 1992.
It is, for all practical purposes, the constitutional standard.

\subsection{Premise 3: Political Self-Interest Is the Source of the Problem}

The diagnosis of gerrymandering as an institutional-integrity failure---
map-drawing captured by those with a personal stake in the outcome---is not
a partisan position.
It is the position of both parties when they are the victim.

Republicans challenged the Maryland congressional map in 2011, which was
designed to flip a Republican-held district to Democratic.
The case eventually reached the Supreme Court as part of the consolidated
redistricting cases decided alongside \rucho{}.\footnote{%
  \textit{Rucho v.\ Common Cause} consolidated two cases: a Republican-drawn
  map in North Carolina challenged by Democrats, and a Democratic-drawn map
  in Maryland challenged by Republicans.
  The Court treated both as presenting the same issue.
  \rucho{}, 588 U.S. at 690--692.
}
Democrats challenged the Wisconsin legislative map, the North Carolina
congressional map, and the Texas congressional map.
Republicans challenged the Maryland congressional map, the New York
congressional map, and the Illinois legislative map.

Each challenge used similar language: the map was drawn to entrench the
drawing party's power, not to serve constituents.
Each challenge was, at its core, an argument that politicians should not
draw their own districts.

This bipartisan consensus on the diagnosis has produced bipartisan action
on the remedy.
Independent redistricting commissions, which remove map-drawing from
legislators and assign it to a panel of citizens or experts, have been
adopted in Arizona, California, Colorado, Michigan, and thirteen other states.
Most of these reforms were passed by ballot initiative, and most commanded
bipartisan support from voters even where they were opposed by legislators
of both parties.

The \dia{} takes the independent-commission insight one step further.
Rather than substituting human commissioners for human legislators---
with all the selection, deliberation, and potential for capture that
commissions entail---it substitutes a mathematical algorithm whose
parameters are entirely determined by federal census data.
No commission needed.
No deliberation required.
The census tells the algorithm everything it needs to know.

\subsection{Premise 4: The Interstate and Intrastate Problems Are Isomorphic}

This is the logical hinge on which the argument turns.
If the interstate and intrastate problems are structurally different, the
analogy fails and the argument collapses.
We argue that they are structurally identical.

\paragraph{The interstate problem.}
Given: 50 states, each with a population $P_s$.
Total seats: $K = 435$.
Minimum: each state gets at least one seat.
Goal: assign $k_s$ seats to each state such that $\sum k_s = K$ and
the average district population $P_s / k_s$ is as equal across states
as possible, minimizing the percentage deviation $|P_s/k_s - P_t/k_t|$
for all state pairs $(s, t)$.
Method: Huntington-Hill.

\paragraph{The intrastate problem.}
Given: a state with $m$ geographic units (census tracts), each with a
population $p_i$.
Total districts: $k$ (the state's apportioned seat count).
Minimum: each district gets at least one unit.
Goal: assign each unit to one of $k$ districts such that the district
populations are as equal as possible.
Method: ???

The mathematical objects in the two problems are different---states with
populations versus census tracts with populations, 435 seats versus
$k$ districts---but the \emph{structure} of the optimization is identical.
Both problems require:
\begin{enumerate}
  \item Partitioning a set of weighted units into exactly $k$ groups.
  \item Minimizing the population deviation across groups.
  \item Respecting a minimum-allocation constraint (one unit per group).
\end{enumerate}

The only structural difference is the adjacency constraint.
In the interstate problem, any set of states can be combined into a
``district'': there is no geographic requirement that congressional
states be contiguous (and indeed, they are not, since Alaska and Hawaii
are not contiguous with the mainland).
In the intrastate problem, districts must be geographically contiguous:
a district whose tracts are not connected is constitutionally suspect.\footnote{%
  Contiguity is universally required by state redistricting statutes,
  though it is not explicitly mandated by the federal constitution.
  Every state that has addressed the question requires contiguity,
  and no federal court has suggested that non-contiguous districts
  would be permissible.
}

\paragraph{The isomorphism and its limits.}
The adjacency constraint distinguishes the intrastate problem in an important respect:
the interstate apportionment problem assigns seats to states (which are fixed, pre-existing
entities), while the intrastate redistricting problem must find contiguous regions
(which are not fixed---they must be constructed from the adjacency graph).
The Huntington-Hill priority rule, when extended to the intrastate setting, therefore
applies to the \emph{bisection tree structure}---determining the sequence and proportions
of recursive splits---rather than to the identification of specific census tracts for
each region.
The isomorphism is in the decision procedure (priority sequence $\to$ factorization $\to$ tree),
not in the feasible set (states vs.\ contiguous connected subgraphs).
Formally, \ar{} applies the \hh{} priority rule to choose how to partition the state at each
level of the bisection tree, and delegates to METIS the task of finding the minimum-edge-cut
connected subgraph partition consistent with that choice.
This division of labor---\hh{} for the tree structure, METIS for the contiguous realization---is
what makes the intrastate extension computable.
The result is a doctrinal extension grounded in the same equal-proportions principle,
applied at the next level down in the apportionment hierarchy.

\subsection{The Logical Step: If HH Is Right at Level~$n$, It Is Right at Level~$n+1$}

The argument does not require any new constitutional theory.
It requires only the recognition that the same principle---equal proportions,
implemented via Huntington-Hill---applies at each level of the apportionment
hierarchy.

Level~1: Allocate 435 seats among 50 states.
Method: Huntington-Hill.
Authority: 2~U.S.C.~\usca{}~2a(a), upheld in \textit{Department of Commerce
v.\ Montana}.

Level~2: Allocate $k$ districts among $m$ census tracts within a state.
Method: Huntington-Hill (via \ar{}).
Authority: the same principle, applied at the next level down.

This is not a novel judicial construction.
Courts routinely apply the same principle at multiple levels of a hierarchy.
The Equal Protection Clause applies at the federal, state, and local levels.
The First Amendment applies to federal, state, and local governments.
The principle of one person, one vote (\wesberry{}) already applies at the
state level.
Applying the equal-proportions method (\hh{}) at the intrastate level
is the natural extension of a principle the Court has already applied
at the interstate level.

The Congress that enacted 2~U.S.C.~\usca{}~2a(a) in 1941 did not have the
computational tools to implement \hh{} at the intrastate level.
We do now.
The \dia{} corrects an 85-year gap in the statutory implementation of the
apportionment principle.

\subsection{The One-Sentence Law}

The conclusion of the four-step argument is a single sentence of federal
statutory law.
We propose it as 2~U.S.C.~\usca{}~2a(c) [new]:

\begin{statute}
\noindent\textbf{Proposed 2 U.S.C.\ \usca{}\ 2a(c) [new].}
Each State shall apportion its congressional districts by applying the
Huntington-Hill priority rule to its census tracts, using as input only
the prime factorization of the State's apportioned seat count and the
federal TIGER/Line adjacency graph for the census year in which the
apportionment was performed, with a seed determined by the formula
\texttt{SHA-256(census\_release\_id \textbar\textbar\ `DIA\_SEED\_V1')}.
\end{statute}

Every word in this sentence is load-bearing.

\begin{itemize}
  \item \textit{Huntington-Hill priority rule}: specifies the mathematical
    method, preventing practitioner substitution.
  \item \textit{census tracts}: specifies the geographic unit, preventing
    gerrymandering by unit selection.
  \item \textit{only the prime factorization of the State's apportioned
    seat count}: limits inputs to the factorization of $k$, excluding
    all political data.
  \item \textit{federal TIGER/Line adjacency graph}: specifies the
    geographic source, preventing manipulation of the underlying graph.
  \item \textit{census year in which the apportionment was performed}:
    ties the geographic data to the apportionment cycle, preventing
    mid-decade redrawing using updated data.
  \item \textit{SHA-256(census\_release\_id \textbar\textbar\
    `DIA\_SEED\_V1')}: makes the random seed content-derived and publicly
    computable, eliminating seed manipulation as a vector for gaming.
\end{itemize}

The full implementation details---the VRA overlay, the county-preservation
default, the gcd-incompatibility fallback---are addressed in
Section~\ref{sec:implementation}.
But the operative principle is captured in these 67 words.
Everything else is implementation detail, governed by the principle they embody.
